% Chi porta cosa nel DC-link: la corrente del ponte e' a gradini, quella della batteria e'
% continua, il condensatore porta la differenza. Forma qualitativa su tre periodi di commutazione.
\begin{tikzpicture}
  \pgfplotsset{dcax/.style={width=0.82\textwidth, height=3.6cm, xmin=0, xmax=3,
    xtick={0,1,2,3}, grid=major, grid style={black!10},
    tick label style={font=\scriptsize}, label style={font=\footnotesize},
    legend style={font=\scriptsize, draw=none, fill=white, fill opacity=0.85, text opacity=1,
                  at={(1.0,1.0)}, anchor=north east, legend columns=2}}}
  \begin{axis}[dcax, name=sopra, ymin=-5, ymax=45, ytick={0,20,40}, xticklabels={},
               ylabel={current [A]}]
    % zero vettori: nessuna corrente; vettori attivi: la corrente di una fase o di due
    \addplot[farancio, thick, const plot] coordinates {
      (0,0) (0.15,36) (0.35,10) (0.65,36) (0.85,0)
      (1,0) (1.15,36) (1.35,10) (1.65,36) (1.85,0)
      (2,0) (2.15,36) (2.35,10) (2.65,36) (2.85,0) (3,0)};
    \addlegendentry{bridge input $i_\mathrm{in}$}
    \addplot[fblu, very thick, domain=0:3] {17.4};
    \addlegendentry{battery $I_\mathrm{dc}$, its average}
  \end{axis}
  \begin{axis}[dcax, at={(sopra.south)}, anchor=north, yshift=-3mm, ymin=-25, ymax=25,
               ytick={-20,0,20}, xlabel={time [switching periods]}, ylabel={current [A]}]
    \addplot[fgrigio, thick, const plot] coordinates {
      (0,-17.4) (0.15,18.6) (0.35,-7.4) (0.65,18.6) (0.85,-17.4)
      (1,-17.4) (1.15,18.6) (1.35,-7.4) (1.65,18.6) (1.85,-17.4)
      (2,-17.4) (2.15,18.6) (2.35,-7.4) (2.65,18.6) (2.85,-17.4) (3,-17.4)};
    \addlegendentry{capacitor $i_C = i_\mathrm{in} - I_\mathrm{dc}$, zero average}
  \end{axis}
\end{tikzpicture}
